% CHƯƠNG 3 - Rubric: Mô hình và Phương pháp/giải pháp/thuật toán
\chapter{Phương pháp Future-Guided Learning}
\label{chuong:phuongphap}

Chương này trình bày chi tiết phương pháp FGL: công thức chưng cất tri thức dùng để kết nối
teacher và student (mục~\ref{sec:fgl-loss-detail}), cách phương pháp hiện thực hóa khái niệm
``tương lai'' khác nhau trên từng bài toán cụ thể (mục~\ref{sec:cochetuonglai} - đây là phần
phân tích quan trọng nhất của chương), kiến trúc các mô hình được dùng
(mục~\ref{sec:kientruc}) và quy trình huấn luyện tổng thể (mục~\ref{sec:quytrinh}). Nội dung
được đối chiếu trực tiếp với mã nguồn công khai của tác giả (kho
\texttt{SkyeGunasekaran/FutureGuidedLearning}, commit \texttt{5fa24ea}) để làm rõ các chi tiết
mà bài báo không nêu tường minh.

\section{Công thức chưng cất tri thức trong FGL}
\label{sec:fgl-loss-detail}

Như đã nêu ở Chương~\ref{chuong:baitoan}, hàm mất mát của student là tổ hợp của một hàm mất
mát chuẩn so với nhãn thật và một hàm mất mát chưng cất so với teacher:
\begin{equation}
\begin{split}
	\mathcal{L}_{\text{student}} = &\ \alpha \cdot \mathcal{L}_{\text{CE}}\!\left(y^{(S)}, f_S(\mathbf{x})\right) \\
	&+ (1-\alpha) \cdot T^2 \cdot D_{\mathrm{KL}}\!\left(\operatorname{softmax}_T\!\left(f_T(\mathbf{x})\right) \,\middle\|\, \operatorname{softmax}_T\!\left(f_S(\mathbf{x})\right)\right).
\end{split}
	\label{eq:fgl-loss-ch3}
\end{equation}
Hệ số $T^2$ nhân vào thành phần KL là một kỹ thuật chuẩn trong chưng cất tri thức
\cite{hinton2015distillation}: khi $T$ tăng, độ lớn gradient của softmax có nhiệt độ giảm theo
tỉ lệ $1/T^2$, nên nhân lại $T^2$ giúp giữ thang gradient ổn định khi thay đổi $T$, không phải
một đóng góp riêng của bài báo này. Tham số $\alpha \in [0,1]$ quyết định mức độ tin vào nhãn
thật so với tin vào teacher: $\alpha = 1$ loại bỏ hoàn toàn ảnh hưởng của teacher (tương đương
mô hình baseline không dùng FGL); $\alpha = 0$ khiến student học hoàn toàn từ teacher, không
dùng nhãn thật.

Đối chiếu với mã nguồn, hai cài đặt của bài báo trên hai bài toán viết công thức này theo hai
cách khác nhau: tệp \texttt{FGL\_CHBMIT.py} nhân hệ số $(1-\alpha)$ ở bên ngoài tích
$D_{\mathrm{KL}} \cdot T^2$, còn hàm \texttt{KL()} trong \texttt{mackey\_glass/utils/utils.py}
trả về giá trị đã nhân sẵn $(1-\alpha)$ bên trong. Hai cách viết này về toán học là
\textit{tương đương}, chỉ khác ở chỗ đặt dấu ngoặc.

\section{Cơ chế hiện thực hóa ``tương lai'' theo từng bài toán}
\label{sec:cochetuonglai}

Công thức~\eqref{eq:fgl-loss-ch3} chỉ xác định \textit{cách kết nối} teacher và student;
phần quyết định ``tương lai'' nghĩa là gì nằm ở cách teacher và student được \textit{huấn
luyện} với dữ liệu và nhãn khác nhau. Đây là phần linh hoạt nhất của phương pháp, được hiện
thực hóa theo hai cách khác nhau tùy loại bài toán.

\subsection{Mackey-Glass: lệch cửa sổ thời gian đầu vào}

Với bài toán hồi quy trên chuỗi Mackey-Glass, không có khái niệm ``nhãn khác nhau'' giữa hai
mô hình như bài toán phân loại; thay vào đó, FGL làm lệch \textit{cửa sổ thời gian đầu vào}.
Cho một chuỗi $\mathbf{y} = (y_1, \dots, y_T)$, với mỗi vị trí $i$, nhãn tại mốc tuyệt đối
$t = i + L + H - 1$ được lấy ra, trong đó $L$ là độ dài cửa sổ nhìn lại (lookback window) và
$H$ là tầm dự báo (forecasting horizon). Teacher và student dùng cùng công thức lấy nhãn này
nhưng với hai cặp $(L, H, \text{offset})$ khác nhau, trình bày trong
Bảng~\ref{bang:mg-offset}.

\begin{table}[H]{Cách lệch cửa sổ giữa teacher và student trên Mackey-Glass}
	\label{bang:mg-offset}
	\centering
	\begin{tabular}{lccc}
		\toprule
		\textbf{Mô hình} & \textbf{Lookback window} & \textbf{Tầm dự báo} & \textbf{Offset đầu vào} \\
		\midrule
		Teacher & $L$ (giống student) & $1$ (chỉ 1 bước) & $H - 1$ \\
		Student & $L$ & $H$ (tầm dự báo thật) & $0$ \\
		\bottomrule
	\end{tabular}
\end{table}

Vì offset của teacher bằng $H-1$, teacher xuất phát từ mốc $i + (H-1)$ và chỉ cần ngoại suy
$1$ bước, nhưng mốc đích cuối cùng vẫn là $t = \left(i+(H-1)\right) + L + 1 - 1 = i+L+H-1$,
\textit{trùng khớp} với mốc đích của student. Nói cách khác, teacher và student cùng nhắm tới
một điểm đích tuyệt đối trên chuỗi, nhưng teacher xuất phát từ một điểm \textit{gần đích hơn}
student, nên bài toán của teacher dễ hơn nhiều (chỉ ngoại suy 1 bước thay vì $H$ bước), trong
khi vẫn có thể dùng để chưng cất tri thức về cùng một đích dự báo cho student.

Vì chưng cất tri thức theo công thức~\eqref{eq:fgl-loss-ch3} yêu cầu đầu ra của teacher và
student là phân phối xác suất rời rạc, bài toán hồi quy được chuyển thành bài toán phân loại
bằng cách \textit{rời rạc hóa}: khoảng giá trị quan sát được của chuỗi được chia thành
\texttt{num\_bins} lớp bằng hàm phân khoảng (tương tự \texttt{np.digitize}), biến mỗi giá trị
liên tục thành chỉ số của lớp chứa nó. Đây là một lựa chọn thiết kế hợp lý nhưng cũng là một
giới hạn: độ phân giải của dự báo cuối cùng bị chặn bởi số lượng lớp \texttt{num\_bins} được
chọn.

\subsection{CHB-MIT: lệch định nghĩa nhãn}

Với bài toán dự báo cơn động kinh trên CHB-MIT, teacher và student không lệch cửa sổ thời
gian mà được huấn luyện trên \textit{hai định nghĩa nhãn khác nhau} của cùng loại tín hiệu EEG,
tương ứng với hai bài toán khác nhau, trình bày trong Bảng~\ref{bang:chbmit-nhan}.

\begin{table}[H]{Hai bài toán khác nhau cho teacher và student trên CHB-MIT}
	\label{bang:chbmit-nhan}
	\centering
	\small
	\begin{tabular}{p{2.6cm}p{5.6cm}p{5.6cm}}
		\toprule
		 & \textbf{Teacher} & \textbf{Student} \\
		\midrule
		Tên bài toán & Seizure Detection (phát hiện cơn \textbf{đang} xảy ra) & Seizure Prediction (dự báo cơn \textbf{sẽ} xảy ra) \\
		Nhãn ictal lấy từ & đoạn tín hiệu \textit{trong lúc} co giật & đoạn tín hiệu \textbf{SOP} = 30 phút, kết thúc cách thời điểm co giật \textbf{SPH} = 5 phút \\
		Độ khó & Dễ - thấy trực tiếp biến cố & Khó - chỉ thấy tín hiệu trước biến cố, chưa có dấu hiệu rõ ràng \\
		\bottomrule
	\end{tabular}
\end{table}

Teacher đóng vai trò ``người đứng ở tương lai'', nhìn thấy trực tiếp tín hiệu trong lúc cơn
động kinh xảy ra, nên bài toán của teacher dễ và đạt độ chính xác tự nhiên cao hơn. Student
đóng vai trò ``người đứng ở quá khứ'', chỉ thấy tín hiệu \textbf{trước} cơn động kinh (giai
đoạn tiền triệu, preictal), phải dự báo trước khi biến cố thực sự xảy ra. Khi chưng cất tri
thức, đầu vào preictal/interictal của student được đưa qua \textit{cả} teacher và student;
teacher chưa từng được huấn luyện trực tiếp trên dữ liệu preictal, nhưng các đặc trưng tần số
nó học được để nhận ra cơn động kinh (từ bài toán detection) vẫn mang thông tin hữu ích để
``cảm nhận'' tín hiệu preictal giống cơn động kinh tới mức nào, từ đó dẫn dắt student.

\subsection{AES: teacher dùng chung theo loài (ngoài phạm vi tái hiện của báo cáo)}

Ngoài Mackey-Glass và CHB-MIT, bài báo gốc còn thực nghiệm trên bộ dữ liệu điện não trong sọ
(intracranial EEG) AES (American Epilepsy Society Seizure Prediction Challenge)
\cite{brinkmann2016aes}, gồm 5 chó (\texttt{Dog\_1}-\texttt{Dog\_5}) và 2 người
(\texttt{Patient\_1}, \texttt{Patient\_2}). Mục này chỉ trình bày cơ chế FGL trên bộ dữ liệu
này ở mức lý thuyết, dựa trên việc đọc mã nguồn và tài liệu của tác giả, để thấy được tính tổng
quát của phương pháp; báo cáo \textbf{không tái hiện thực nghiệm} trên bộ dữ liệu AES (lý do
nêu ở Chương~\ref{chuong:dulieu}).

AES cũng là bài toán dự báo cơn động kinh như CHB-MIT, nhưng thiếu nhãn ictal đầy đủ cho từng
đối tượng riêng lẻ để tự huấn luyện một teacher riêng theo từng bệnh nhân/vật. Tác giả giải
quyết bằng cách huấn luyện \textit{hai teacher dùng chung} - một cho nhóm chó và một cho
nhóm người - trên dữ liệu ictal/interictal gộp chung của cả nhóm, rồi tái sử dụng cùng một
teacher đó khi chưng cất tri thức cho mọi đối tượng cùng loài. Đây là một cách khác để hiện
thực hóa ``tương lai'': tri thức mà teacher nắm được không đặc thù theo cá thể, mà là tri thức
\textit{gộp chung của cả loài} về đặc trưng tần số của cơn động kinh.

\section{Kiến trúc các mô hình}
\label{sec:kientruc}

\subsection{Mackey-Glass: RNN}

Teacher, baseline và student trên Mackey-Glass đều dùng cùng kiến trúc: một mạng RNN 2 lớp
(128 đơn vị ẩn) nối tiếp hai tầng kết nối đầy đủ (fully-connected), với tầng cuối có số đầu ra
bằng \texttt{num\_bins} (số lớp sau khi rời rạc hóa). Việc dùng cùng một họ kiến trúc cho cả
ba mô hình giữ cho khác biệt về hiệu năng chỉ phản ánh tác dụng của cơ chế FGL, không bị nhiễu
bởi khác biệt kiến trúc.

\subsection{CHB-MIT và AES: CNN-LSTM và MViT}

Hai bộ dữ liệu điện não dùng chung hai lựa chọn kiến trúc, định nghĩa trong
\texttt{models/model.py}:

\begin{itemize}
	\item \textbf{CNN\_LSTM\_Model}: ba tầng tích chập 3 chiều (Conv3D, xử lý spectrogram đa
	kênh như một ``ảnh 3D'' thời gian-tần số-kênh) kèm các tầng gộp (pooling), sau đó làm
	phẳng (flatten) và đưa qua một mạng LSTM 3 lớp (512 đơn vị ẩn), cuối cùng là hai tầng kết
	nối đầy đủ cho ra 2 đầu ra (ictal/interictal). Đây là kiến trúc mặc định cho cả teacher và
	student/baseline.
	\item \textbf{MViT} (``multi-channel ViT''): mỗi kênh EEG được embed thành một chuỗi token
	riêng, đưa qua một bộ mã hóa Transformer \textit{riêng cho từng kênh} (không chia sẻ trọng
	số giữa các kênh), lấy token \texttt{[CLS]} của mỗi kênh rồi nối lại, qua một tầng kết nối
	đầy đủ để phân loại. Ý tưởng là coi mỗi kênh EEG như một ``modal'' riêng, học đặc trưng
	không gian-thời gian độc lập trước khi kết hợp (fuse) thông tin giữa các kênh.
\end{itemize}


\section{Quy trình huấn luyện tổng thể}
\label{sec:quytrinh}

Hình~\ref{hinh:quytrinh-fgl} tóm tắt quy trình ba giai đoạn (stage) mà cả hai bài toán đều
tuân theo, chỉ khác nhau ở cách định nghĩa dữ liệu/nhãn của từng giai đoạn như đã trình bày ở
mục~\ref{sec:cochetuonglai}.

\begin{figure}[H]{Ba giai đoạn huấn luyện của Future-Guided Learning}
	\label{hinh:quytrinh-fgl}
	\centering
	\begin{tikzpicture}[
		node distance=0.5cm,
		hop/.style={draw, rounded corners=2pt, align=center, text width=8cm, minimum height=0.95cm, font=\small},
		tieude/.style={font=\small\bfseries},
		mui/.style={-{Stealth[length=2mm]}, thick}
	]
		\node[hop] (s1) {\textbf{Giai đoạn 1 - Teacher}\\Huấn luyện teacher trên nhiệm vụ dễ (tương lai gần hơn), độc lập, không liên quan đến student};
		\node[hop, below=of s1] (s2) {\textbf{Giai đoạn 2 - Baseline}\\Huấn luyện một mô hình trên nhiệm vụ khó của student, không dùng teacher (đối chứng)};
		\node[hop, below=of s2] (s3) {\textbf{Giai đoạn 3 - Student (FGL)}\\Huấn luyện student trên nhiệm vụ khó, có teacher (đã đóng băng) dẫn dắt qua hàm mất mát \eqref{eq:fgl-loss-ch3}};
		\node[hop, below=of s3] (s4) {\textbf{Đánh giá}\\So sánh Baseline với Student trên cùng tập kiểm tra để đo tác dụng của FGL};

		\draw[mui] (s1) - (s2);
		\draw[mui] (s2) - (s3);
		\draw[mui] (s3) - (s4);
		\draw[mui, dashed] (s1.east) - ++(0.7cm,0) |- (s3.east);
	\end{tikzpicture}
\end{figure}

Mũi tên nét đứt thể hiện việc teacher, sau khi đã huấn luyện xong ở Giai đoạn 1, được dùng lại
(ở trạng thái đóng băng, không huấn luyện tiếp) để tính hàm mất mát chưng cất trong Giai đoạn
3. Baseline ở Giai đoạn 2 có vai trò là mô hình đối chứng: cùng kiến trúc, cùng dữ liệu huấn
luyện với student, nhưng không có teacher dẫn dắt ($\alpha = 1$ trong công thức
\eqref{eq:fgl-loss-ch3}), dùng để đo \textit{riêng} tác dụng của cơ chế chưng cất tri thức,
tách biệt khỏi tác dụng của kiến trúc hay dữ liệu.
